% Exposé IV. The lens: tangent and cotangent methods
% Sources: theory/framework.md (Exposé IV), site/sections/07-lens.html, theory/propositions.json (IV.2-IV.4).
\section{The lens: tangent and cotangent methods}\label{sec:lens}

Backpropagation pairs every map with its reverse derivative, so a method can shape what the model expresses on the
forward side or compress the gradient on the backward side. For linear methods and optimizers that read only their
gradients (no weight decay, no clipping of the full gradient), the two sides give the same weights, and GaLore becomes
one-sided LoRA with a moving frame and carried optimizer state.

\subsubsection{Two sides of one port}\label{sec:lens-hook}

A layer does two jobs in training. Forward, it turns parameters and inputs into outputs. Backward, it turns a cotangent
at the output into cotangents at the input and at the parameters. Backpropagation keeps the pair together, and a
category theorist calls it a \emph{lens} (\bgref{B.17}).

Every method in Exposés~\ref{sec:slice} to~\ref{sec:fibres} changes the forward job, replacing the parameter port
$\Theta$ by a smaller space $Q$ and a map $\rho:Q\to\Theta$. A second family edits only the backward job. GaLore
\citep{zhao2024galore} compresses the weight gradient $G$ to $P^\top G$, with $P$ from an SVD of $G$, runs Adam on that
small matrix and maps the result back with $P$. Every weight trains; only the optimizer state is small. LISA
\citep{pan2024lisa} and MeZO \citep{malladi2023fine} follow the same pattern.

When $\rho$ is linear the two families meet. The pullback of the forward method is then the compression of a backward
method, and an optimizer that reads only its gradients cannot tell the two apart. The equivalence between GaLore and
one-sided LoRA is due to \citet{xtorrobahennigen2025duality}, who prove a duality between linear gradient
transformations and adapters (their Thm~1 and Cor.~3, with weight decay in their App.~B). The frozen-$A$ case under
SGD, with random projections, appears earlier in \citet{hao2024flora}. Our additions are the lens bookkeeping, the
explicit hypotheses on the optimizer, and the table that sorts the cotangent family
(Table~\ref{tab:lens-schedule}).

\subsubsection{Backpropagation as a lens}\label{sec:lens-functor}

In the language of \thmref{Definition}{I.1}, a layer is a parametrised map $f:P\times X\to Y$. Its reverse derivative
$\mathsf R[f](x,y)=J_f(x)^\top y$ takes an output cotangent to an input cotangent, parameters included, and
backpropagation is the 2-functor
\begin{equation}\label{eq:lens-functor}
  \Para(\mathsf R):\Para(\mathcal C)\longrightarrow\Para(\Lens(\mathcal C)),\qquad f\longmapsto(f,\mathsf R[f]).
\end{equation}
The parametrised-map construction and the reading of backpropagation as a functor into lenses are due to
\citet{fong2017backprop} and \citet{cruttwell2021categorical} (\bgref{B.16}, \bgref{B.17}).

A reparametrisation $\rho$ is a 2-cell, so the functor turns it into a lens $(\rho,\mathsf R[\rho])$, and a
reparametrisation therefore acts on both sides of the parameter port. Its tangent side is the image, which is what the
method can \emph{express} (Exposé~\ref{sec:image}). Its cotangent side $\mathsf R[\rho](q,\gamma)=D\rho_q^\top\gamma$
pulls the weight gradient back to $Q$ and, with the rates $\Lambda$, induces the cometric $K=D\rho\,\Lambda\,D\rho^\top$
(\thmref{Definition}{III.1}, \bgref{D.23}), which is how the method \emph{moves}. A lens need not come from a map,
though, and a backward half written down directly gives the second family. The two kinds of lens on the parameter port
are
\[
\begin{tikzcd}[column sep=large,row sep=normal]
  Q \arrow[r,"\rho"] & \Theta \\
  T^*_qQ & T^*_{\rho(q)}\Theta \arrow[l,"D\rho_q^{\top}"']
\end{tikzcd}
\qquad\qquad
\begin{tikzcd}[column sep=large,row sep=normal]
  & \Theta \arrow[r,"\id_\Theta"] & \Theta \\
  T_\theta\Theta & S \arrow[l,"a_t"'] & T^*_\theta\Theta \arrow[l,"c_t"']
\end{tikzcd}
\]
on the left a forward method, whose backward half is forced by its forward half, and on the right a backward method,
whose forward half is the identity and whose backward half is chosen freely.

\begin{definition}{IV.1}{forward and backward methods}
\begin{itemize}
  \item A \textbf{forward} method acts on the parameter port by a lens in the image of $\mathsf R$, namely
  $(\rho,\mathsf R[\rho])$.
  \item A \textbf{backward}, or cotangent-side, method acts by a lens whose forward part is $\id_\Theta$ and whose
  backward part is a \emph{compression} $c_t:T^*_\theta\Theta\to S$ into a state space $S$, paired with an
  \emph{anchor} $a_t:S\to T_\theta\Theta$. The update is
  \begin{equation}\label{eq:lens-backward-update}
    \theta\leftarrow\theta-\eta\,a_t\big(\mathrm{opt}(c_t(\nabla L))\big).
  \end{equation}
\end{itemize}
Table~\ref{tab:lens-backward} lists examples.
\end{definition}

\begin{table}[t]
\centering\footnotesize
\caption{Backward methods as compression--anchor pairs (\thmref{Definition}{IV.1}). The last row is the only reading
of Flora as published in the form of \thmref{Definition}{IV.1}, the trivial lens with all compression inside
$\mathrm{opt}$.}\label{tab:lens-backward}
\begin{tabularx}{\linewidth}{@{}>{\raggedright\arraybackslash}p{0.27\linewidth}>{\raggedright\arraybackslash}X>{\raggedright\arraybackslash}p{0.24\linewidth}@{}}
\toprule
\headrow \textbf{Method} & \textbf{Compression $c_t$} & \textbf{Anchor $a_t$} \\
\midrule
GaLore \citep{zhao2024galore} & $P_t^\top$, with $P_t$ from a gradient SVD & $P_t$ \\
GaLore with a random projector \citep{xtorrobahennigen2025duality,hao2024flora} & as GaLore, with $P_t$ random & $P_t$ \\
LISA \citep{pan2024lisa} & layer-block coordinate projection & the same projection \\
MeZO \citep{malladi2023fine} & $g\mapsto\langle z_t,g\rangle$, estimated by finite differences without a backward pass &
$\lambda\mapsto\lambda z_t$ \\
\midrule
Flora as published \citep{hao2024flora} & identity: only the momentum or accumulation buffer is compressed, by a random
$A_t$, inside $\mathrm{opt}$ & identity (full-space Adafactor step) \\
\bottomrule
\end{tabularx}
\end{table}

A backward method has its own cometric $K_t=a_tc_t$ (\thmref{Definition}{III.1}), for instance $K(G)=PP^\top G$ for
GaLore under SGD. The update must lie in the image of the anchor. Read with their projections as anchors, Flora as
published \citep{hao2024flora}, APOLLO \citep{zhu2024apollo} and Fira \citep{chen2024fira} make updates that leave
$\operatorname{im}a$, so they are not backward methods in the sense of \thmref{Definition}{IV.1} beyond the trivial
reading with compression and anchor the identity, which the last row of Table~\ref{tab:lens-backward} records for
Flora. Flora compresses its momentum or gradient accumulator with a
random projection and hands the decompressed tensor to Adafactor \citep{shazeer2018adafactor}, which acts in the full
space, so its update need not lie in the image of the projection. On resampling it transports its momentum to the new
frame. It fits \thmref{Definition}{IV.1} with a nontrivial anchor only in an idealised form with an SGD or heavy-ball
base.

\subsubsection{The degree-one collapse}\label{sec:lens-collapse}

If $\rho(C)=W_k+\mathcal LC$, training $C$ and running the backward method with $c=\mathcal L^*$, $a=\mathcal L$ feed
the optimizer the same gradients $\mathcal L^*G$ and so make the same weight updates. A curved $\rho$, such as LoRA's
$BA$, separates them after one step, and so does an optimizer that reads the parameter tensor, which is $C$ in one
procedure and $W$ in the other.

\begin{proposition}{IV.2}{the degree-one collapse; known, \citet{xtorrobahennigen2025duality} for GaLore (their Thm~1
and Cor.~3, App.~B for weight decay) and \citet{hao2024flora} for random projections}
Let $\rho_k(C)=W_k+\mathcal LC$ with $\mathcal L$ linear and $C_0=0$. Let $\mathrm{opt}$ be a first-order optimizer,
with no coupled or decoupled weight decay, whose state and output depend only on the sequence of gradients it is fed,
and which does not read the parameter tensor: SGD, heavy-ball or Nesterov momentum, Adam or AMSGrad, or Adafactor with
\texttt{scale\_\allowbreak parameter=False} (fixed or relative step, with its factorisation applied to the compressed gradient in
the same tensor shape as $C$). Run both procedures below with the same implementation of $\mathrm{opt}$, on the same
minibatches, and from the same optimizer state, fresh or carried over. Then they produce the same $W$-trajectory over a
period:
\begin{itemize}
  \item training $C$ with $\mathrm{opt}$ and merging at the end of the period;
  \item the backward method with $c=\mathcal L^*$ and $a=\mathcal L$, started at $W_k$.
\end{itemize}
AdamW with $\lambda>0$, LAMB and LARS trust ratios, parameter-scaled Adafactor (\texttt{scale\_\allowbreak parameter=True}, the
default in HF and in \texttt{galore\_torch}'s \texttt{GaLoreAdafactor}) and global-norm clipping of the full gradient
(such as HF Trainer's default \texttt{max\_grad\_norm=1.0}, whenever it triggers) all break the identity. Injectivity
of $\mathcal L$ is not needed.

In particular, on each matrix that GaLore projects (its other parameters train in full in both procedures), GaLore with
a projector $P_k$ fixed on each period equals one-sided LoRA $W_k+P_kC$ (frozen $B=P_k$, or frozen $A$ when projecting
on the right), merged and re-based every period, with $C$'s optimizer state (moments and step counter) carried over
unchanged across re-basings, with no reset. $P_k$ must be exactly GaLore's matrix, signs and column order included; its
span alone does not suffice, because the state is reused without transport. ReLoRA-style resets give a different
trajectory. The identity is exact with weight decay $0$, no full-gradient clipping, and an optimizer from the admissible
class; \texttt{galore\_torch}'s \texttt{GaLoreAdafactor} defaults to \texttt{scale\_\allowbreak parameter=True} and is excluded.
\texttt{GaLoreAdamW} is the transformers-style AdamW, which adds $\varepsilon$ to $\sqrt v$ before the bias
correction, so the LoRA side must use the same implementation. GaLore's scale $\alpha$ multiplies the optimizer output,
so it acts as a learning-rate factor; a LoRA-style scale inside $\rho$ gives a different trajectory under Adam with
$\varepsilon>0$.
\end{proposition}

\begin{remark*}{scale and implementation}
Moving GaLore's $\alpha$ inside $\rho$ as a LoRA scale multiplies the compressed gradient by $\alpha$. Under Adam this
gives GaLore's trajectory only at $\varepsilon=0$, or with the LoRA side's $\varepsilon$ multiplied by $\alpha$ as well.
This is the exchange of scale for learning rate and $\varepsilon$ on the hyperparameter torus of
\thmref{Theorem}{III.10}. Both sides must also run the same implementation: \texttt{galore\_torch}'s AdamW adds
$\varepsilon$ to $\sqrt v$ before the bias correction, as in transformers, and swapping in \texttt{torch.optim.Adam} on
the LoRA side moves the trajectory measurably (see the check below).
\end{remark*}

\noindent\labelsf{Num}\enspace In a toy check under Adam the two procedures agree to $7\times10^{-16}$ over one
period, and to $5\times10^{-15}$ over four periods with the state carried on both sides. Against the official
\texttt{galore\_torch} optimizers (least squares, $W\in\R^{6\times10}$, $r=2$, $T=7$, four periods, total movement
$0.046$), the two procedures agree to $8\times10^{-17}$ with the state carried over. Resetting the LoRA side's state
separates them by $0.029$, weight decay $0.1$ by $0.020$, HF-default Adafactor by $0.030$ (against $6\times10^{-17}$
with \texttt{scale\_\allowbreak parameter=False}), and flipping the sign of one column of $P_k$, with the state carried, by $0.12$.
Using \texttt{torch.optim.Adam} on the LoRA side moves the trajectory by about $4\times10^{-5}$, and placing $\alpha$
inside $\rho$ by $0.019$ at $\varepsilon=10^{-6}$, with exact agreement at $\varepsilon=0$. The companion site reruns a
check of this kind live (Appendix~\ref{app:repro}).

\begin{explains}
Beyond one-sided LoRA, GaLore adds two ingredients: a frame read off the gradient's SVD, and the carry-over of optimizer
state across re-basings. The carry-over involves no change of basis. Moments written in the coordinates of $P_{k-1}$
are read as coordinates for $P_k$. Projection-aware state updates, such as LDAdam \citep{xrobert2025ldadam}, and Flora
transport the state instead. The equivalence classifies the cotangent methods whose optimizer satisfies its hypotheses.
Table~\ref{tab:lens-schedule} sorts them by where the projector comes from and how often it changes; its LoRA and
ReLoRA row is a contrast row outside the equivalence.
\end{explains}

\begin{table}[t]
\centering\footnotesize
\caption{The cotangent family by projector (rows) and schedule (columns). \thmref{Proposition}{IV.2} covers the first
three rows under its optimizer hypotheses, except EVA and LoRA-GA, which initialise two-sided LoRA and sit in the frozen
column only for their first step. The last row is a contrast row outside the equivalence.}\label{tab:lens-schedule}
\begin{tabularx}{\linewidth}{@{}>{\raggedright\arraybackslash}p{0.19\linewidth}>{\raggedright\arraybackslash}X>{\raggedright\arraybackslash}X>{\raggedright\arraybackslash}X@{}}
\toprule
\headrow \textbf{Projector / schedule} & \textbf{Frozen} & \textbf{Re-based periodically} & \textbf{Resampled every step} \\
\midrule
random & LoRA-FA \citep{zhang2023lora}, with $K\propto GA_0^\top A_0$ under SGD$^{a}$ & GaLore with a random projector
\citep{xtorrobahennigen2025duality,hao2024flora}; idealised Flora$^{b}$ & MeZO \citep{malladi2023fine} (rank $1$,
global), up to its finite-difference estimator \\
\addlinespace
gradient or data SVD & EVA \citep{paischer2024parameter}, LoRA-GA \citep{wang2024lorab} (first step only) & GaLore
\citep{zhao2024galore} & GaLore with $T=1$ \\
\addlinespace
coordinate blocks & BitFit \citep{benzaken2021bitfit}, FISH \citep{sung2021training}, LN-tuning
\citep{zhao2023tuning} & LISA \citep{pan2024lisa} & block coordinate descent \\
\addlinespace
two-sided, state-dependent (contrast) & LoRA \citep{hu2021lora} & ReLoRA \citep{lialin2023relora} & none in use \\
\bottomrule
\end{tabularx}
\par\smallskip
\begin{minipage}{\linewidth}\scriptsize
$^{a}$\,For the original method under SGD. HF PEFT's LoRA-FA optimizer preconditions by $(A_0A_0^\top+\delta I)^{-1}$,
which in the undamped limit gives the row-space projector $GP_V$ of \thmref{Proposition}{III.4}(c); its Adam layer
makes the update nonlinear in $G$.\quad
$^{b}$\,Flora's Alg.~2 with an SGD or heavy-ball base, with the momentum transported at resampling instead of carried
over; Flora as published, with an Adafactor base, is outside \thmref{Definition}{IV.1} and outside
\thmref{Proposition}{IV.2}.
\end{minipage}
\end{table}

\thmref{Proposition}{IV.2} covers the first three rows of Table~\ref{tab:lens-schedule} under its optimizer hypotheses,
except EVA and LoRA-GA, which initialise two-sided LoRA and sit in the frozen column only for their first step; it
covers MeZO up to the finite-difference error of its estimator. LoRA's map $(B,A)\mapsto W_0+sBA$ has degree
two, so its cometric moves with the state (\thmref{Observation}{III.2}) and the collapse does not apply. The other
frozen entries gain nothing from merge-and-restart (\thmref{Theorem}{V.3}, corollary~2); their re-based neighbours move
the frame and can reach all of weight space. The LoRA-FA cometric holds for the original method of
\citet{zhang2023lora} under SGD; the LoRA-FA optimizer of HF PEFT \citep{xmangrulkar2022peft} gives the row-space projector of
\thmref{Proposition}{III.4}(c), as noted under the table.

Flora proper is not in the table. Its Adafactor step acts on the full gradient, with parameter scaling by default, so
its steps are not confined to $\operatorname{im}P_t$ and \thmref{Proposition}{IV.2} does not apply; like LDAdam, it
transports its optimizer state. Only the idealised form with an SGD or heavy-ball base sits in the random, re-based
cell, and there its momentum is transported rather than carried.

\subsubsection{When a backward method is a forward one}\label{sec:lens-integrability}

A backward method chooses at each weight a space of directions $D_\theta=\operatorname{im}a(\theta)$. It reads as a
forward method only if these directions fit together into $k$-dimensional surfaces, the images of immersive
parametrisations (\bgref{C.4}). Frobenius' theorem tests this (\bgref{D.25}, \auxref{B.31}). The bracket $[X,Y]$, the gap by which a
small loop along $X$, $Y$, $-X$, $-Y$ fails to close, must stay in $D$. Discrete steps and the optimizer's metric ask
for more.

\begin{proposition}{IV.3}{integrability; after Frobenius' theorem}
Let $D$ be a smooth distribution of constant rank $k$ on an open set $U\subseteq\Theta$, for instance the anchor
distribution $D_\theta=\operatorname{im}a(\theta)$ of a time-independent backward method. Then $D$ is involutive on
$U$ if and only if through every $\theta\in U$ passes an injective immersion $\rho_\theta:(\R^k,0)\to(U,\theta)$ with
$\operatorname{im}D\rho_\theta=D$ along its image. A single such $\rho$ through $\theta_0$ gives involutivity only
along $\operatorname{im}\rho$. For example, $D=\spanop\{\partial_x,\ \partial_y+xz\,\partial_z\}$ on $\R^3$ has
$[\partial_x,\partial_y+xz\,\partial_z]=z\,\partial_z\notin D$ off $z=0$, yet the plane $z=0$ is an integral manifold.

The dynamic content needs separate hypotheses.
\begin{itemize}
  \item In continuous time, if $D$ is involutive, a backward flow with $\dot\theta\in D_\theta$ stays on the leaf
  through $\theta_0$.
  \item In the discrete update \eqref{eq:lens-backward-update} of \thmref{Definition}{IV.1}, the iterates stay on a
  leaf for all step sizes and all optimizer outputs $s$ (all losses) if and only if that leaf contains
  $(\theta+D_\theta)\cap U$ at every iterate, that is, if and only if the leaves are affine; an anchor constant over
  the period gives this (\thmref{Proposition}{IV.2}).
  \item Under gradient flow, dynamic equivalence with a forward lens $(\rho,\mathsf R[\rho])$ forces matching frames,
  $a_tc_t=D\rho\,\Lambda\,D\rho^\top$ along the trajectory. Hence $K=ac$ is symmetric positive semidefinite and, for
  injective $a$, $c=Ma^\top$ with $M$ symmetric positive semidefinite on $S$, possibly depending on $\theta$. If in
  addition $\rho$ is an immersion with $\dim Q=\rk D$, then $\rho$ is a local isometry from $(Q,\Lambda^{-1})$ onto the
  leaf with metric $(aMa^\top)^{-1}|_D$, so that metric must be flat. Forward methods with gauge fibres escape this
  requirement: $\rho(B)=BB^\top$ on symmetric matrices induces the curved Bures--Wasserstein metric, and balanced
  full-rank LoRA also induces a curved leaf metric. For discrete steps, exact agreement for every step size further
  forces $\rho$ to be affine along each horizontal line $q+t\,\Lambda D\rho_q^\top g$.
\end{itemize}

\emph{Instances.} Every projector-type backward method of Table~\ref{tab:lens-backward} (GaLore with any period $T$
and any projector, its random-projector variant, LISA, MeZO) has an anchor that does not depend on $\theta$ within its
period. Its distribution is constant, hence trivially involutive with affine leaves (vacuously at $T=1$, where a period
is a single step). Under the optimizer hypotheses of \thmref{Proposition}{IV.2} (no weight decay of either kind, no
full-gradient clipping; GaLore and MeZO need weight decay $0$, and LISA's AdamW with positive decay violates them),
these methods are rebasing schemes of linear forward methods with period $T$, including $T=1$, whose reach
\thmref{Theorem}{V.3} describes. Each MeZO step is a step of the linear method $\theta_t+\lambda z_t$, up to its SPSA
estimator, and GaLore with $T=1$ is one-sided LoRA re-based every step. What separates these methods from LoRA-type
methods is the resampling schedule together with the carry-over of optimizer state; integrability plays no part in it.
Flora as published (Adafactor base, with a full-rank current-gradient term in the code), APOLLO (a
low-rank-estimated scaling applied to the full gradient) and Fira (a scaled full-rank residual) are not of the form of
\thmref{Definition}{IV.1} with their projections as anchors, because their updates leave $\operatorname{im}a$.

Non-involutivity proper needs a $\theta$-dependent anchor. Full-batch GaLore with $T=1$, read as the distribution
$D_W=\operatorname{Hom}(\R^n,U_r(\nabla L(W)))$ (or $\operatorname{Hom}(V_r(\nabla L(W)),\R^m)$ when $m\ge n$, as
GaLore projects) wherever a singular-value gap exists, is involutive if and only if $U_r(\nabla L(W))$ (respectively
$V_r(\nabla L(W))$) is invariant along $D$. That fails generically for every $r\ge1$ once both matrix dimensions are at
least $2$. For every $r$ with $1\le r<\min(m,n)$ it fails for $L=\tfrac12\lVert W\rVert_F^2$, and numerically it fails
for least squares and cross-entropy with several samples. It holds, for instance, for single-sample least squares, and
automatically only when the distribution is a line field ($nr=1$). Minibatch GaLore with $T=1$ and MeZO violate
time-independence, because their anchors depend on step-wise randomness.
\end{proposition}

\noindent\labelsf{Num}\enspace For full-batch GaLore with $T=1$, the relative normal part of the bracket is $0.08$ to
$0.94$ at $r=1$ and $(m,n)\in\{(4,3),(5,4)\}$, for least-squares, tanh-regression and cross-entropy losses with several
samples (an independent replication), and $0.21$ to $0.95$ at $(m,n,r)=(4,3,2)$. It vanishes only for $n=1$ or for
single-sample least squares. At $W=\diag(2,1)$ with $L=\tfrac12\lVert W\rVert_F^2$ and the frame $X_i=u(W)e_i^\top$,
the bracket is $[X_1,X_2]=-\tfrac13e_2e_1^\top\notin D$, and the general diagonal formula behind it was checked by
central differences for several $(m,n,r)$. The cometric of $B\mapsto BB^\top$ takes the same value at $B$ and at $BR$,
$R\in\Ort(2)$, to $10^{-12}$. The companion site recomputes the bracket live at $W\in\R^{3\times4}$
(Appendix~\ref{app:repro}).

\begin{explains}
The lens 2-functor \eqref{eq:lens-functor} separates tangent-side from cotangent-side methods. Frobenius integrability
is the test for when a $\theta$-dependent anchor foliates weight space by images of forward methods. For the PEFT zoo it
does no discriminating work. Every cotangent method whose update lies in $\operatorname{im}a$, whose optimizer does not
read the parameter, and which has no decoupled decay on $W$ and no clipping of the full gradient (GaLore with any
projector and MeZO at weight decay $0$, LISA without decay) is, period by period, a forward linear method on a moving
base (Exposé~\ref{sec:base}). At $T=1$ this reading is automatic.
\end{explains}

\begin{analogy*}{the lens duality}
``Tangent side'' and ``cotangent side'' are exact words here. A forward method acts on tangent vectors through
$D\rho$, and a backward method acts on cotangent vectors through $c$. Calling GaLore and one-sided LoRA \emph{dual} is
exact only at degree one, and only for optimizers that do not read the parameter tensor, with no weight decay. Beyond
that range the word organises the zoo and proves nothing.
\end{analogy*}

\subsubsection{What the backward pass must remember}\label{sec:lens-cone}

Fine-tuning memory consists of weights, gradients and optimizer state, and the activations stored for the backward
pass. Shrinking $Q$ shrinks the gradients and the optimizer state. As for the activations, a frozen linear box's reverse
derivative $W^\top y$ ignores its input, and cotangents are needed only downstream of something trainable, so only a
cone through the network costs activation memory.

\begin{observation}{IV.4}{backward reach and the backprop cone; an upper bound}
Assume reverse mode in which each box keeps its own residual, with no checkpointing across boxes, and take boxes at the
granularity where adapters are separate boxes (for LoRA-FA, $A_0$ and $B$ separately). A \emph{VJP residual} is an
input, an output, a sufficient statistic such as a ReLU mask, or a recipe to recompute one of these inside the box.
Then it suffices to keep a VJP residual for
\begin{itemize}
  \item each trainable box whose parameter-VJP depends on its input: linear and multiplicative boxes such as $W$, $A$,
  $B$ and (IA)$^3$ scales; additive boxes such as biases and prompts need nothing;
  \item each \emph{nonlinear} box, bilinear boxes such as the attention products included, on a directed path from a
  trainable box to the loss.
\end{itemize}
Deterministic frozen linear boxes, and all boxes off this cone, need nothing; stochastic linear boxes on the cone
(dropout) need their mask or random-number state. Alternatively, the $k$ coordinates of $P^\top\nabla L$ can be
computed exactly by $k$ forward-mode JVPs, with no reverse pass and $O(1)$-layer activation memory, each JVP costing a
small constant multiple (about $1.5$ to $3$) of a forward pass. This always uses less activation memory than
un-checkpointed reverse mode, but it is time-competitive with one VJP (about $3$ forward passes) only for $k$ of about
$1$ or $2$. It also needs $P$ before the passes, so it cannot produce GaLore's SVD refresh, where $P$ depends on
$\nabla L$, more cheaply than a backward pass ($\nabla L$ itself would take $\lvert\theta\rvert$ JVPs).
\end{observation}

\begin{explains}
The bound reads off the activation memory of the zoo.
\begin{itemize}
  \item \textbf{LST} \citep{sung2022lst}. The side network reads the backbone through trainable down-projection
  ladders, so the backbone's backward pass never runs.
  \item \textbf{LoRA-FA} \citep{zhang2023lora}. With \texttt{lora\_dropout=0} (the HF default) it stores $A_0x$ ($r$
  numbers per token) instead of $x$ ($n$ numbers); with dropout on, the $n$-element mask is stored too. The saving is
  realised only where no neighbouring nonlinear box already keeps $x$.
  \item \textbf{Top-layer adaptation} truncates the cone, under the no-checkpointing assumption. With HF's reentrant
  gradient checkpointing, PEFT calls \texttt{enable\_input\_require\_grads()}, which makes the embedding outputs require
  gradients, so cotangents run down to the embeddings and top-layer adaptation no longer truncates the cone.
  \item \textbf{LISA} \citep{pan2024lisa} trains the embedding and the LM head at every step, so its cone is the whole
  network and it gets no truncation. On the cone, the unsampled layers' frozen linear boxes drop their input residuals
  (about $20\%$ of a block's activations at LLaMA proportions without FlashAttention), and relative to LoRA it also
  drops the adapter residuals, in line with the LISA paper's report of less activation memory than LoRA. Most of its
  saving against full fine-tuning is in gradients and optimizer state. A variant that froze the embedding would
  truncate the cone at its lowest sampled layer.
  \item \textbf{Prompt tuning} \citep{lester2021power} has a $Q$ as small as that of VeRA \citep{kopiczko2023vera},
  FourierFT \citep{gao2024parameter} or (IA)$^3$ \citep{liu2022few}, and the maximal cone, the whole network, which it
  shares with every method that trains something in the first layer. Its prompt tokens also enlarge every stored
  activation.
  \item \textbf{MeZO} \citep{malladi2023fine} is the case $k=1$, with the JVP estimated by a central difference and no
  reverse pass. Its estimator is unbiased for the gradient of the Gaussian-smoothed loss and, for a loss with Lipschitz
  Hessian, has bias $O(\varepsilon^2)$ against $\nabla L$; ReLU networks lack a Lipschitz Hessian. Exact unbiasedness,
  $\mathbb E[(\nabla L\cdot z)z]=\nabla L$, holds for the forward gradient computed with an exact JVP. GaLore, its
  random-projector variant, Flora and LISA all take the reverse route.
\end{itemize}
\end{explains}

Ladder side-tuning is due to \citet{sung2022lst}, prompt tuning to \citet{lester2021power}, and the zeroth-order
estimator in this setting to \citet{malladi2023fine}. That reverse derivatives of linear maps ignore their inputs is the
observation behind many activation-saving designs; the lens of \citet{cruttwell2021categorical} makes it visible box by
box.

So every cotangent method meeting the hypotheses of \thmref{Proposition}{IV.2} is, period by period, a linear forward
method on a moving base. What it reaches, and what a quantiser does to the base, are the subject of
Exposé~\ref{sec:base}.
